% Generated draft: requires mathematical and source review.
\documentclass[11pt]{article}
\usepackage[margin=25mm]{geometry}
\usepackage{fontspec}
\setmainfont{Latin Modern Roman}
\usepackage{amsmath,amssymb,mathtools,unicode-math}
\setmathfont{Latin Modern Math}
\usepackage{enumitem,needspace,xurl,hyperref}
\setlength{\parindent}{0pt}
\setlength{\parskip}{5pt}
\setlength{\emergencystretch}{3em}
\newcommand{\N}{\mathbb{N}}
\newcommand{\Z}{\mathbb{Z}}
\newcommand{\Q}{\mathbb{Q}}
\newcommand{\R}{\mathbb{R}}
\newcommand{\C}{\mathbb{C}}
\newcommand{\F}{\mathbb{F}}
\newcommand{\A}{\mathbb{A}}
\newcommand{\PP}{\mathbb{P}}
\newcommand{\E}{\mathbb{E}}
\newcommand{\eps}{\varepsilon}
\newcommand{\dd}{\,\mathrm{d}}
\begin{document}
\section*{UnsolvedMath ID: 30003310}
\textbf{Source problem code:} catalogue entry
\subsection{Definitions and mathematical statement}
Let \((R,\mathfrak m)\) be a local ring of finite Krull dimension \(d\), let \(I\subseteq R\) be an ideal, and let \((I_n)_{n\geq 1}\) be a sequence of ideals satisfying \(I_{n+1}\subseteq I_n\). Interpret cofinality with the powers of \(I\) as the mutual containment condition
\[
(\forall k\geq 1)(\exists n\geq 1)\ I_n\subseteq I^k,
\qquad
(\forall n\geq 1)(\exists k\geq 1)\ I^k\subseteq I_n.
\]
Write \(H^i_{\mathfrak m}(M)\) for the local cohomology of an \(R\)-module \(M\) with support in \(\mathfrak m\), and \(\operatorname{length}_R(M)\) for its composition length. For a fixed nonnegative local-cohomology index \(i\), suppose that
\[
(\exists N\geq 1)(\forall n\geq N)\qquad
\ell_n^i:=\operatorname{length}_R H^i_{\mathfrak m}(R/I_n)<\infty.
\]
The supplied question asks: what asymptotic behavior, as \(n\to\infty\), can be established for \(\ell_n^i\) under these hypotheses?

The formulation remains qualified: the record does not explicitly specify commutativity or Noetherianity of \(R\), the intended range of \(i\), or a convention for local cohomology if rings outside the usual commutative Noetherian setting are intended. These conventions require confirmation. The finiteness hypothesis above concerns the chosen index \(i\); simultaneous finiteness for all indices is not stated. No multiplicative compatibility or quantitative comparison between \(I_n\) and powers of \(I\) is supplied. The record specifies neither a normalization nor a particular asymptotic assertion, such as a growth bound or existence of a limit. Consequently, this is a broad asymptotic question rather than a specified conjectural formula.
\subsection{Short English statement}
For a decreasing sequence of ideals cofinal with the powers of an ideal \(I\), study the eventual growth of the finite lengths \(\operatorname{length}_R H^i_{\mathfrak m}(R/I_n)\) at a fixed index \(i\). The ring conventions, index range, and intended asymptotic target require clarification. The supplied status is \texttt{partially\_solved}; this draft does not independently verify that assessment.
\subsection{Sources}
\begin{itemize}
\item Local UnsolvedMath registry: supplied complete record for permanent ID \(30003310\), titled \emph{Asymptotic Lengths of Local Cohomology Modules}, in the Oberwolfach Reports Open Problems collection; exact source problem code \texttt{catalogue entry}. No individual proposer is supplied. The recorded proposed year is 2016. The supplied \texttt{corrected\_verified} label is qualified by a verification note describing only notation, typography, and TeX-delimiter cleanup while preserving hypotheses, quantifiers, and subquestions; it is not independent mathematical verification here.
\item Supplied report citation: \emph{Asymptotic Phenomena in Local Algebra and Singularity Theory} (2017), Oberwolfach Reports, pp.\ 3199--3259, DOI 10.4171/owr/2016/57. Source URL and literature-source label ``original OWR report'': \url{https://doi.org/10.4171/owr/2016/57}. The background identifies Oberwolfach Reports as documenting workshops and their open-problem discussions.
\item Supplied literature source, labeled ``Dao--Montaño'': \url{https://arxiv.org/abs/1705.05033}. The record supplies no further bibliographic details for this link.
\item Literature assessment dated exactly \texttt{2026-08-21}: supplied dated curation, not independently verified. It labels the problem partially solved and attributes to Dao and Montaño asymptotic upper bounds and actual rational limits for important classes, including suitable monomial ideals, while stating that the broad local-ring question remains unsettled. The background characterizes this as dated literature triage derived from the supplied OWR curation and calls for independent verification of time-sensitive status and cited claims. Neither linked source has been independently consulted for this draft.
\end{itemize}
\end{document}
